%%%%%
%
\section{Longueur de Monin-Obukhov (formule approximative)}

\itemstyle{Op\'erateur~:} {\sf Mol\_approx}
\index{Mol\_approx}
\newline

%%%%%OBJET%%%%
\itemstyle{Objet~:}
Longueur de Monin-Obukhov calcul\'ee par la formule approximative � partir des covariances et de la temp\'erature moyenne.
LMO est la hauteur (entre -200 et +200 m) � laquelle la production thermique de turbulence l'emporte sur la production m�canique :
\begin{center}
L>0 : stabilit� ; L=0 : convection libre ; L<0 : instabilit� 
\end{center}


%%%%%ARGUMENTS%%%%
\bigskip
%%%%%ENTREES%%%%
\begin{tabarg}
\multicolumn{3}{l}{\titretab {Entr\'ees}}\\
\hline
$COV_{wt}$ & Vecteur, \texttt{float} & Covariance $\overline{W'T'}$ (\textit{mK/s})\\
$COV_{wu}$ & Vecteur, \texttt{float} & Covariance $\overline{W'U'}$ (\textit{$m^2/s^2$})\\
$COV_{wv}$ & Vecteur, \texttt{float} & Covariance $\overline{W'V'}$ (\textit{$m^2/s^2$})\\
$T_{airmoy}$ & Vecteur, \texttt{float} & Temp\'erature de l'air moyenne (\textit{C} ou \textit{K})\\
$t_{int}$ & Coeff., \texttt{float} & Temps d'int\'egration pour le calcul (\textit{s})\\
\end{tabarg}

\medskip
%%%%%SORTIES%%%%
\begin{tabarg}
\multicolumn{3}{l}{\titretab {Sorties}}\\
\hline
$L_{mon}$ & Vecteur, \texttt{float} & Longueur de Monin Obukhov (\textit{m})\\
\end{tabarg}


%%%%%FORMULE%%%%
\itemstyle{Formule}

%%Mol_approx

Le r\'esultat $L_{mon}$ est calcul\'e � partir des covariances. La base de temps doit \^etre commune.

\begin{align*}
u_* &= (COV_{wv}^{2}+COV_{wv}^{2})^{\frac{1}{4}}\\
\intertext{On en d\'eduit :}
L_{mon} &= \frac{-u_*^3.T_{airmoy}}{k.g.COV_{wt}}
\end{align*}

avec \newline
$k$ : constante de Kolmogorov choisie �gale � 0.4\\
$g$ : constante de gravit� �gale � 9.81 $ms^{-2}$





%%%%%AUTEUR%%%%
\itemstyle{Auteur~:}
CNRM/GMEI/TRAMM

%%%%%REFERENCE%%%%
\itemstyle{R\'ef\'erences~:}
